\documentclass[11pt]{article}
\usepackage[a4paper,margin=1in]{geometry}
\usepackage{array}
\usepackage{url}

\setlength{\parindent}{0pt}
\setlength{\parskip}{0.75em}

\begin{document}

\begin{center}
{\Large Assignment 1B Learning Reflection}\\[0.4em]
{\large Cryptography and Cryptanalysis Upskilling for the SEPAR Thesis}\\[0.8em]
Noah Ostle -- 24740589
\end{center}

\section*{Introduction}

This reflection follow Gibb's reflective cycle, and aims to evaluate how I executed the independent learning plan that I set for my research project on the SEPAR lightweight IoT cipher. My goal was not only to increase general cryptography knowledge, but to develop the specific capability needed to make progress on a cryptanalytic thesis: reading a cipher design critically, identifying which parts of its construction might expose weaknesses, and the skills necessary to perform both mathematical and structural analysis.

The reflection focuses on three external learning activities from my plan: the TU Graz Cryptography lectures, the TU Graz Cryptanalysis lectures, and the EC-Council Certified Encryption Specialist material. These activities were chosen because they directly support my thesis methodology. My current thesis draft - with the help of the knowledge gained from, these 3 units - identifies several preliminary observations that still need proof or strong empirical validation, but that seem promising for the upcoming semester of work on my thesis. The learning activities were therefore evaluated against a practical test: whether they improved my ability to understand and investigate those observations.

\section*{1. Description}

In the original plan, I proposed to complete a structured set of external activities alongside thesis work. The three activities reflected on here totalled 56 hours and 5.6 learning-plan credit points (excluding the UTS module);

\textbf{TU Graz Cryptography WS Lectures.} I planned 20 hours, worth 2.0 credit points in my learning plan. The official TU Graz/ISEC course page describes Cryptography as a Masters level course in the Cryptology and Privacy area, delivered as lectures plus practical classes. The 2025/26 public schedule lists lectures and practicals, with the full current lecture recordings uploaded and publicly available on the internet. The content covers the design and mathematical evaluation of authenticated encryption, hashing, symmetric primitives such as AES-GCM and SHA-3, asymmetric encryption and signatures, key-exchange protocols, privacy-preserving protocols, and introductory cryptanalysis for the purpose of secure design. In practice, I used the online lecture material asynchronously during the semester. I watched the lectures and worked through the examples alongside the tutor, but I did not have a formal exam result or certificate from this activity. This fact however, should not undermine the contents extreme usefulness and rigor; despite the main mode of content being non interactive lectures, they proved fundamental to my cryptographic knowledge, and subsequently my thesis, as I will detail later.

\textbf{TU Graz Cryptanalysis SS Lectures.} I planned 16 hours, worth 1.6 credit points. The official TU Graz/ISEC course describes Cryptanalysis as a masters level course focused on modern cryptanalytic attacks and how they guide cryptographic design. The 2026 schedule lists lectures lectures and practicals, with the full current lecture recordings uploaded and publicly available on the internet. The course content includes classic and quantum algorithms for factoring and discrete logarithms, block cipher cryptanalysis, differential, linear and algebraic attacks, hashing function and stream cipher cryptanalysis, lattices, continued fractions, implementation aspects, and exercises in which students implement attacks against cryptanalytic challenges. For my own plan, I again used the online lecture material and worked examples. This meant I could complete the technical content as self-study rather than as a formally enrolled student with graded assignments, similar to the above module.

\textbf{ECES Certified Encryption Specialist.} I planned 20 hours, worth 2.0 credit points. EC-Council describes ECES as a 3-day, 20-hour program delivered through in-person, self-study or live-online modes, with six modules and a 50-question, two-hour multiple-choice exam under exam code 212-81. The public outline covers classical cryptography, symmetric cryptography and hashes, number theory and asymmetric cryptography, applications such as PKI, TLS, VPNs and disk encryption, cryptanalysis, and quantum computing. My actual completion differed from my original expectation of taking the paid certification immediately; I purchased the official study pack in textbook form, read through the full material, and paused at each chapter to test retention through Quizlet flashcards. As a further self-assessment artefact, I completed a CertHero ECES 212-81 practice-question set and scored 95 percent. This is evidence of engagement and recall, but I unfortunately have not had sufficient time (between work, uni, and other commitments) to actually schedule and complete the exam.


\section*{2. Feelings}

With the support of my supervisor, this thesis started out more as a learning activity than an application of knowledge - I admitted that although I had a good basis in mathematics and cybersecurity already, I had not undertaken much rigorous and dedicated cryptographical study, and that I wanted to use this project as a chance to do so. I could identify \textit{some} exmaples of suspicious behaviour in SEPAR, but I did not yet have enough background knowledge or technical skill to prove whether an observation was cryptographically meaningful, statistically valid/useful, or a false positive. This was uncomfortable because cryptanalysis requires above all else, rigor; claiming a result without full confidence is one way to have your research rejected.

The TU Graz material initially made that gap more obvious. The Cryptography lectures showed how much design rationale sits behind even apparently simple primitives, while the Cryptanalysis lectures showed that attacks are not just a list of named techniques. They require a model, an observable, a distinguisher or recovery goal, complexity estimates, and a careful explanation of what assumptions are being made. That made my early SEPAR observations feel less complete, but it also made the path forward clearer.

By contrast, the ECES material felt more familiar and structured. It was broader, less research-specialised, and more industry focussed. That made it easier to measure progress chapter by chapter. Using Quizlet flashcards and practice questions gave me a visible feedback loop that the TU Graz courses did not provide. However, I also became aware that scoring well on applied cryptography questions is different from being able to produce an original cryptanalytic argument. The confidence it gave me was useful, but it had to be kept in proportion given what I came to understand was its focus on general applied cryptographic knowledge, rather than the intricate low level understanding required for research in the field.

\section*{3. Evaluation}

The strongest part of the plan was its alignment with my research. The TU Graz Cryptography lectures helped me understand why a lightweight cipher cannot be evaluated only by whether it uses recognisable components such as S-boxes, rotations, XOR, modular addition, an IV, or a large nominal key. I became more aware that security depends on how those components interact. That directly improved my ability to discuss SEPAR's hybrid design, its small block size, and its stateful stage structure without treating any one feature as automatically secure or automatically broken.

The Cryptanalysis lectures were the most technically valuable activity. They gave me a clearer framework for moving from "this looks unusual" to "this is the adversarial model, this is the observable, this is the statistic or structure being tested, and this is the claim the experiment can actually support." In particular, the linear, differential and algebraic material mapped directly onto my planned hypotheses about SEPAR. It also helped me see that my thesis should not promise a completed break before the individual weakness points are proven or rejected.

The main weakness in my execution was the lack of formal mastery artefacts for the TU Graz courses (although there is little my plan could have done to address this shy of enrolling at the university). I watched the lectures and worked through the examples, but because I was not completing the courses as an enrolled student, I do not have graded submissions, exam transcripts, or certificates. I addressed this by making the thesis itself the evidence of learning;
although perhaps less visible to an outside observer, the quality and accuracy of the research I have already produced is far and away of a higher quality than what I would have been capable of at the start of the semester.

The ECES activity was successful in a different way. Reading the official study pack forced me to consolidate the broader cryptographic context: modes, IVs, key management, hashes, asymmetric primitives, PKI, TLS, VPN's, disk encryption, steganography, common implementation mistakes, and introductory cryptanalysis. The 95 percent practice result I achieved on previous practice exams available online suggests that I retained the material at a factual level. It's limitation is that it suffers (in my application at least) from being "a mile wide, but an inch deep". This has some benefits, as it exposed me to many aspects of cryptography that I might not have considered if only approaching it as a purely mathematical and academic endeavour - the business need or privacy requirement is what necessitates cryptographic primitives to have certain mathematical features, and is what provides the security context and basis for the adversarial model when examining a cipher. That being said, it helped me explain operational implications, and worked side by side with the more technical and research focussed TU Graz courses.

\section*{4. Analysis}

\subsection*{TU Graz Cryptography: understanding design choices}

The Cryptography lectures changed how I read SEPAR's design. Before this learning activity, I was more likely to assess a cipher by listing its components: block size, key size, IV size, S-boxes, rotations, and number of stages. After the lectures, I became more focused on the design goals and security assumptions behind those components. For example, the TU Graz course content explicitly connects modern primitives to confidentiality, integrity, authentication, hashing, authenticated encryption, protocol use, and cryptanalysis. That perspective matters for SEPAR because the cipher is not a standard block cipher or a standard stream cipher. It is described as a hybrid construction, so I need to be precise about what security property is being claimed and under what model.

This learning directly contributed to my current thesis work. In the introduction and methodology draft, I describe SEPAR as a composition of translated stage permutations rather than as a single opaque algorithm. That model lets me ask whether the public 16 bit permutation for a fixed IV behaves like an independently keyed random permutation, or whether the eight stage structure leaks information. The Cryptography lectures helped me see why this distinction is important. A small domain is not automatically broken, infact the course goes over several similar smaller domain algorithms that are considered secure! What it does do however, is make complete codebook analysis feasible. The research question is therefore not simply "is 16 bits too small?", but whether the small domain interacts badly with the stage decomposition, state translations, or key partition.

The lightweight cryptography and symmetric primitive focussed parts of the course also helped me position SEPAR against other designs. My draft now uses examples such as Hummingbird-1, PRESENT, GIFT, KeeLoq, and KTANTAN to make a more careful point: lightweight design choices can be valid, but they need a security argument that matches the construction. This changed the tone of my writing. Instead of saying SEPAR is weak because it is \textit{too} lightweight, I now frame the thesis as an investigation into whether SEPAR's particular efficiency choices leave exploitable structure.

\subsection*{TU Graz Cryptanalysis: turning observations into tests}

The Cryptanalysis lectures were most important for the technical method. Differential cryptanalysis taught me that "difference" is not a generic concept. In SEPAR, XOR differences, modular additive differences, and differences before or after a state translation may behave differently. This directly shaped the differential section of my methodology. If a table has an unknown output translation, modular additive differences may cancel that translation in a way that XOR differences do not. That observation gives me a concrete reason to test additive difference scores rather than treating all differential experiments as equivalent, and is something I would not have even known to consider prior to viewing the lectures.

The linear cryptanalysis material helped me understand why null hypotheses and sample sizes matter. My original plan already referred to detecting linear bias, but the lectures made the requirement more rigorous: a useful bias claim needs a specified statistic, an expected random baseline, a sampling method, and enough trials to distinguish real structure from noise. This connects to SEPAR because preliminary experiments suggest weak diffusion or triangular behaviour in parts of the implementation. The next step cannot be a vague claim that the cipher is biased. It must be a reproducible test that states which bits or quotient views are being measured, what distribution is expected under a random permutation, and how often the behaviour survives through the full cascade. Previously, I would not have known any of these factors which go into making an argument for statistical weakness convincing - they all were familiar mathematically, but the content of the lecture taught me the importance of the methodology, both in presenting findings to others, as well as making sure they are rigorous and not false positives yourself.

The algebraic and block cipher attack material also influenced how I think about stagewise analysis. SEPAR's key is partitioned into eight 32-bit stage keys. The natural cryptanalytic hope is that some observable property identifies or filters one stage at a time, reducing the problem below a 256-bit exhaustive search. The course examples helped me see both the attraction and the risk of this idea. A stage peeling strategy is only meaningful if the score used to recognise a stage remains valid after unknown translations, state updates, and composition with neighbouring stages. This is the entire driving point behind why my thesis draft now separates the stage cascade, quotient structure, triangularity, differentials, and matched contexts into individual weakness points that must be tested independently before being combined, and is very important to my current works accuracy.

The TU Graz exercises were also useful because they modelled a discipline of implementing attacks on small, controlled challenges. That directly applies to my project. SEPAR's 16-bit block size makes full codebook experiments practical, so I can build small reproducible harnesses rather than relying on intuition. The research should proceed by generating codebooks under controlled keys and IVs, testing quotient maps and differential scores, measuring false positives, and then deciding whether any property remains visible from an attackers position.

\subsection*{ECES: connecting technical findings to deployment impact}

ECES contributed most strongly to the practical and professional side of the thesis. The official outline includes symmetric and asymmetric cryptography, hashes, block and stream ciphers, modes of operation, IVs, PKI, TLS, VPN's, disk encryption, common cryptography mistakes, cryptanalysis, and quantum-computing impacts. This breadth helped me place SEPAR in the operational context for which it was proposed: constrained IoT systems.

The most useful ECES lesson for my thesis was that implementation and deployment details can change the meaning of an algorithmic weakness. A theoretical weakness may have little impact if the adversarial model is unrealistic, while a smaller weakness can be serious if it aligns with how a device actually uses keys, IVs, resets, or repeated messages. This influenced my methodology section on adversarial models and public observables. For example, a complete codebook is only relevant if an attacker can obtain reset access, chosen plaintexts, or repeated IV conditions. A matched context experiment is only relevant if chosen prefixes or comparable state conditions are plausible. ECES helped me keep those practical questions in view.

The ECES material also improved how I can communicate the significance of the work. My thesis is not just an internal exercise - it is an assessment of how well I can effectively and scientifically communicate. SEPAR was proposed for IoT environments, so a final result should explain what the weakness, partial weakness, or negative result means for developers and security professionals. The study material on standards, correct and incorrect deployment, and common implementation mistakes gave me language for that discussion, and helped me think about how to translate a cryptanalytic result into advice about whether SEPAR should be trusted, avoided, or further analysed before deployment.

\section*{5. Conclusion}

Overall, I executed the most important parts of the plan effectively. I completed the external learning activities I needed for the thesis, and the learning has changed the way I understand the SEPAR cipher. The main improvement is that I now approach the cipher less like a black-box puzzle and more like a research object: define the construction, define the adversarial model, identify public observables, test individual structural claims, and only then discuss whether those claims combine into an attack.

The plan as a deliverable was not perfect, although the topics I chose themselves could not have provided a greater understanding of the topic at hand. The TU Graz learning would be stronger as evidence if I had formal graded artefacts. However, the absence of formal proof of mastery does not mean the learning had no effect; It's effect is visible in my thesis progress in the move from broad suspicion to specific questions about technical aspects of the cipher. ECES provided a more measurable study outcome through the textbook, flashcards, and 95 percent practice-question result, but its contribution was broader and more contextual.

The most important conclusion is that my knowledge has increased in breadth and specificity as it pertains to the actual research task. I am now better equipped to design a methodology, analyse the cipher, decide what evidence is required, and to write the results into a formal paper.

\section*{6. Action Plan for Next Semester}

Next semester, my priority will be to use the knowledge gained through these learning activities to finish my thesis to a higher level. The main focus will be on moving from broad understanding and preliminary observations toward a coherent research argument, where each technical claim is supported by an appropriate model, method, and proof commensurate with what is expected in the field.

I will begin by consolidating my understanding of SEPAR as a cryptographic construction, ensuring that my thesis rests on a clear explanation of the cipher's design, intended security goals, and relevant assumptions. This will allow the later analysis to be framed in a principled way, rather than as a collection of unsubstantiated observations.

The next stage will be to apply this knowledge through careful research practice. I will use the methodological skills developed during the semester (through the coursework) to analyse and propose results carefully. Where the research produces strong findings, I will work to express their significance clearly. Where it produces negative or partial results, I will treat those outcomes as valuable contributions to understanding the cipher rather than as failures of the project.

I will also continue developing the broader professional perspective gained from the ECES study. This will help me connect the technical conclusions of the thesis to the practical context in which SEPAR was proposed: lightweight cryptography for constrained IoT environments. A final thesis result should not only describe whether a weakness exists, but also explain what that means for confidence in the design, future analysis, and possible real world use.

Overall, my action plan is to use what I have gained this semester to follow through on my current work, and finish off and present my findings by the end of my thesis. The learning completed this semester has given me the conceptual foundation, technical vocabulary, and cryptographic knowledge needed to do that work. The task now is to apply those skills consistently, so that the final thesis presents a thorough and academically useful assessment of SEPAR.


\end{document}
